% Related work. Literature, not results: citations only from references.bib (verified 2026-10-23); a needed reference
% that is not in it is a \TODO{cite: ...} marker, not an entry. OGC's and Do et al.'s content from kaggle/RELATED_WORK_OGC.md
% (literature notes, paraphrased). A sentence that uses this project's results cites its CLAIMS row, as elsewhere.
\section{Related work}
\label{sec:related}

\paragraph{Compressing 3D Gaussian Splatting.}
3D Gaussian Splatting represents a scene by millions of anisotropic Gaussians, each with a position, a covariance, an opacity and a view-dependent colour stored as spherical-harmonic (SH) coefficients, which make up most of each Gaussian's stored floats \citep{kerbl2023gaussian}.
Compression methods for it divide into those that learn a compact representation during training and those that start from a trained model \TODO{cite: a survey of 3DGS compression, e.g. Bagdasarian et al., 3DGS.zip}.
Training-time methods prune Gaussians with learned masks, learn colour codebooks or hash-grid colour fields jointly with the scene, or anchor Gaussians to structured features \TODO{cite: Lee et al. 2024 (Compact 3D Gaussian representation); Navaneet et al. 2024 (CompGS); Girish et al. 2024 (EAGLES); Lu et al. 2024 (Scaffold-GS); Chen et al. 2024 (HAC)}.
Post-training methods prune Gaussians, reduce the SH degree, or vector-quantise (VQ) the parameters of a trained model, usually followed by a short fine-tuning \TODO{cite: Fan et al. 2024 (LightGaussian); Papantonakis et al. 2024 (reducing the memory footprint of 3DGS)}.
This paper works in the post-training setting, and its host is C3DGS \citep{niedermayr2024compressed}.

C3DGS scores each colour and shape parameter by its sensitivity, the magnitude of the gradient of the rendered training images' total energy (the sum of their pixel values) with respect to it, accumulated over the training views, and gives each Gaussian the largest sensitivity among its colour parameters and, separately, among its shape parameters \citep{niedermayr2024compressed}. % C3DGS paper Sec. 4.1, Eq. 3, App. E; code at 2a234af5: compress.py:40-84, 154-156
It removes the Gaussians whose colour sensitivity is zero, clusters the colours of the others into a codebook with sensitivity-weighted $k$-means, and lets the Gaussians whose colour sensitivity lies above a threshold keep their own colour \citep{niedermayr2024compressed}. % Sec. 4.1, Eqs. 5-6, App. E; compression/vq.py:37-60, 78-117, 155, 250-253; arguments/__init__.py:75, 92
It clusters the shape in the same way, as normalised covariance matrices with a separate scale factor per Gaussian, quantises every parameter except the position to 8 bits and stores positions as 16-bit floats, fine-tunes the model with the quantisation in the loop, and stores the arrays, sorted in Morton order, with DEFLATE compression \citep{niedermayr2024compressed}. % Sec. 4.1, App. F, Sec. 4.2, Sec. 4.3; compression/vq.py:194-238; scene/gaussian_model.py:78-88, 182-198, 450-451, 463, 575; compress.py:231; arguments/__init__.py:72, 90; finetune.py:12-24
This paper changes only C3DGS's colour codebook and keeps the rest of its pipeline as released, fine-tuning included (Section~\ref{sec:setup}). % Am17 b

\paragraph{Observation-Gram compaction.}
OGC \citep{pietroszek2026only} weights a Gaussian's SH colour by what the training cameras saw.
For each Gaussian it accumulates a $16\times16$ observation Gram matrix over the SH basis: a sum over the training views of the outer product of the SH basis at the viewing direction, weighted by the Gaussian's compositing weights in that view, so that a colour error in a direction no training camera observed costs nothing \citep{pietroszek2026only}.
In its own pipeline OGC uses the matrix to project the colour onto fewer SH degrees, to allocate coefficients among Gaussians, and as the metric of a generalised Lloyd VQ of the AC coefficients \citep{pietroszek2026only}.
Hosted in C3DGS, it replaces only C3DGS's colour VQ with that Gram-metric Lloyd VQ on all 16 SH rows, keeps C3DGS's pruning, shape VQ, fine-tuning and container, and reports a paired PSNR gain over its reproduction of C3DGS of +0.49 dB before C3DGS's fine-tuning and +0.09 dB after it, on nine Mip-NeRF 360 scenes \citep{pietroszek2026only,barron2022mipnerf360}. % OGC's Table 1, as A4 and A5 state it; A4, A5
Its paper states a ridge of $10^{-6}$ in the VQ's codeword update, while its released code defaults to $10^{-3}$; this paper uses the code's value (Section~\ref{sec:setup}). % RELATED_WORK_OGC.md 2 b; Am17 b; B7; CLAIMS Gap 12
This paper replicates OGC's host comparison under a pre-registered protocol (Section~\ref{sec:replication}), tests whether the lead survives fine-tuning at equal rate (Section~\ref{sec:finetuning}), and dissects which parts of OGC's VQ carry the lead (Section~\ref{sec:dissection}). % Am17; Am19

Do, Chou and Cheung \citep{do2026transforming} also weight SH coefficients by how they are observed, but with one global Gram matrix over all basis functions of all splats, approximated as a product of spatial, directional and colour factors.
Its square root orthonormalises the coefficients before uniform scalar quantisation and entropy coding; the method codes sparse voxel splats rather than 3DGS models, and has no VQ \citep{do2026transforming}.

\paragraph{Second-order and observation-weighted quantisation.}
Weighting a quantiser's error by a second-order model of the loss predates both methods.
In neural-network compression the Hessian of the loss, or its Gauss--Newton approximation, has been used to choose which weights to prune \TODO{cite: LeCun et al. 1990 (Optimal Brain Damage); Hassibi and Stork 1993 (Optimal Brain Surgeon)}, to weight the $k$-means clustering of weights \TODO{cite: Choi et al. 2017 (Towards the limit of network quantization)}, to set per-layer bit widths \TODO{cite: Dong et al. 2019 (HAWQ)}, and to quantise weights one at a time while the remaining weights compensate \TODO{cite: Frantar and Alistarh 2022 (OBQ); Frantar et al. 2023 (GPTQ)}.
Layer-wise quantisers that minimise the error of a layer's outputs on calibration data weight the error by the Gram matrix of the layer's inputs, which is the same construction applied to activations \TODO{cite: Nagel et al. 2020 (AdaRound); Frantar et al. 2023 (GPTQ)}.
For one Gaussian the rendered colour is linear in its SH coefficients, ignoring the clamp of negative colours, so a Gauss--Newton metric of the image error with respect to those coefficients has the form of OGC's observation Gram; the metric used in this paper is a 16-probe estimate of it, equal to OGC's Gram in expectation by construction (Section~\ref{sec:setup}). % Am17 b; A12; RELATED_WORK_OGC.md 2 a, 3
C3DGS's sensitivity is a first-order, scalar counterpart: one weight per Gaussian, the largest gradient magnitude among its colour parameters, so its clustering distance does not distinguish viewing directions. % C3DGS paper Sec. 4.1, Eq. 5; compress.py:154; compression/vq.py:37-60
OGC's scalar variant, the trace of the Gram matrix times the identity, lies between the two, and Section~\ref{sec:gram-trace} tests the full matrix against it. % Am17 b; B1
This project began with a Gauss--Newton-weighted VQ of the SH codebook in gsplat's PNG codec (GN-VQ), which was never compared with OGC on a gate scene and was retired before the replication (Section~\ref{sec:gnvq}). % s17; Am17 a; C5, D6
